\newpage
\section{Non-Volatile Memory Express (NVMe)}
\textbf{Non-Volatile Memory Express (NVMe)} is a high-performance storage interface and communication protocol designed specifically for 
Solid-State Drives (SSDs). It was developed to overcome the limitations of older storage protocols such as Serial ATA (SATA), 
Advanced Host Controller Interface(AHCI) and Serial Attached SCSI (SAS), which were originally designed for Hard Disk Drives (HDDs). 
By utilizing the high-speed Peripheral Component Interconnect Express (PCIe) bus, NVMe enables significantly \underline{faster data transfer rates}, 
\underline{lower latency},and \underline{higher I/O Operations Per Second} (IOPS), allowing SSDs to achieve their full performance potential.


Unlike traditional storage protocols,used in Hard Drives which allows only 1 queues and 32 commands per queue , NVMe is optimized for modern 
flash memory and supports a highly parallel architecture. It can handle up to 65,535 submission queues and 65,535 completion queues, with each 
queue supporting up to 65,536 commands. This enables multiple CPU cores to communicate with the storage device simultaneously, reducing bottlenecks 
and improving system responsiveness during data-intensive operations.

The increasing demand for high-speed storage has been driven by applications such as cloud computing, artificial intelligence (AI), big data analytics, 
virtualization, content creation, and high-performance gaming. These workloads require rapid access to large volumes of data, where traditional storage 
interfaces often become performance bottlenecks. NVMe addresses these challenges by reducing software overhead, minimizing latency, and maximizing the 
bandwidth provided by PCIe, resulting in faster application loading, improved multitasking, and enhanced overall system performance.


Beyond its speed, NVMe offers \textbf{additional advantages} such as \underline{improved energy efficiency}, \ul{scalability}, and support for advanced 
features including \underline{namespace management}, \ul{end-to-end data} \underline{protection}, and \underline{power management}. These capabilities 
make it suitable for a wide range of environments, from consumer laptops and desktop computers to enterprise servers and large-scale data centers. 
Furthermore, technologies such as NVMe over Fabrics (\textbf{NVMe-oF}) extend NVMe's benefits across high-speed networks, enabling remote storage access 
with minimal performance loss.


\textbf{Installation:} An M.2 NVMe SSD is installed directly into the motherboard's M.2 slot and secured with a single screw, requiring no separate data 
or power cables. It communicates through the high-speed PCI Express (PCIe) interface, providing significantly faster data transfer than SATA. A compatible 
motherboard with an M.2 slot and NVMe support is required.
\cite{lexar_nvme}